\section{La reestructuración de la educación: de enseñar a personas a enseñar a agentes}

\subsection{Micro GPT: la esencia en 200 líneas}

El proyecto Micro GPT de Karpathy condensa el entrenamiento de un LLM en 200 líneas de Python (incluidos los comentarios). Es el resultado de diez años de obsesión con «boiling down»: si no se persigue la velocidad y solo se conserva la esencia del algoritmo, el código necesario resulta extremadamente conciso.

\subsection{La educación mediada por agentes}

\begin{importantbox}{De explicarle a las personas a explicarle a los agentes}
«I'm not explaining to people anymore. I'm explaining it to agents.» Karpathy considera que la educación está atravesando un cambio fundamental:
\begin{itemize}
    \item Una vez que el agente entiende el contenido, puede ofrecer una explicación personalizada según el nivel y el idioma de cada estudiante
    \item El papel del educador se transforma en: escribir para el agente un curriculum skill (guión de habilidad didáctica)
    \item La documentación debería pasar de HTML for humans a markdown for agents
\end{itemize}
«The things agents can't do is your job now. The things agents can do, they can probably do better than you.»
\end{importantbox}

\begin{knowledgebox}{Los «few bits» que el ser humano no puede sustituir}
Karpathy intentó que el agente escribiera Micro GPT desde cero, pero fracasó: no logró alcanzar el grado de simplicidad extrema que él pulió durante diez años de obsesión. Sin embargo, el agente comprende por completo el significado de esas 200 líneas de código y puede explicárselas a cualquier persona. El valor del ser humano está en aportar esa contribución creativa de «few bits»; el resto, el agente puede hacerlo mejor.
\end{knowledgebox}

\subsection{Resumen de la sección}

La educación pasa de «personas que enseñan a personas» a «personas que orquestan agentes que enseñan a personas». La contribución central del educador se condensa en el diseño del curriculum y en esas ideas originales de «few bits» que el agente no puede producir por sí mismo.


